% ECONOMICS_PROBLEM: {"id": "J62-441", "title": "Dynastic persistence beyond parents", "jel_code": "J62", "source_id": "OP-0364"}
% !TeX program = xelatex
\documentclass[11pt,a4paper]{article}
\usepackage[margin=23mm]{geometry}
\usepackage{fontspec}
\setmainfont{Latin Modern Roman}
\usepackage{amsmath,amssymb,mathtools}
\usepackage{xcolor,enumitem,booktabs,longtable,array,xurl,hyperref}
\definecolor{muted}{HTML}{53636C}
\hypersetup{colorlinks=true,urlcolor=blue,linkcolor=blue}
\setlength{\parindent}{0pt}
\setlength{\parskip}{5pt}
\newcommand{\R}{\mathbb R}
\newcommand{\E}{\mathbb E}
\newcommand{\N}{\mathbb N}
\newcommand{\ind}{\mathbf 1}
\newcommand{\Var}{\operatorname{Var}}
\newcommand{\Cov}{\operatorname{Cov}}
\newcommand{\supp}{\operatorname{supp}}
\DeclareMathOperator*{\argmin}{arg\,min}
\DeclareMathOperator*{\argmax}{arg\,max}
\newcommand{\entrymeta}[3]{{\sffamily\small\color{muted}\textbf{\texttt{#1}}\quad|\quad #2\par #3\par}}
\newcommand{\Problemref}[1]{\texttt{#1}}
\newcommand{\JELref}[2]{\texttt{#2} (JEL \texttt{#1})}
\begin{document}
\section*{Dynastic persistence beyond parents}
\textbf{Problem ID:} \texttt{J62-441}\quad \textbf{JEL:} \texttt{J62}\par
% BEGIN ECONOMICS STATEMENT
\label{problem:OP-0364}
{\sffamily\small\color{muted}Original subject group: Inequality and economic mobility\par}
\label{definitions:problem:30007646}
\textbf{Audit classification:} Published research agenda. Full paper checked. The original p. 11 locator was inaccurate; mechanism discrimination is an explicit agenda, not a universal persistence law.
\subsection*{Definitions and precise research target}
\textbf{Complete editorial problem.} Consider linked families observed for at least three generations. Let \(Y_{ig}\) be a consistently measured lifetime economic-status variable for family line \(i\) in generation \(g\), and let \(\widetilde Y_{ig}=Y_{ig}+\epsilon_{ig}\) denote the observed proxy with explicitly modeled measurement error \(\epsilon_{ig}\). The descriptive task is to estimate the joint transition law \(P(Y_{ig}\mid Y_{i,g-1},Y_{i,g-2},\ldots,X_i)\), where \(X_i\) includes cohort and institutional context. Test whether more distant ancestry adds predictive information after reliably measured parental status. A nonzero grandparental regression coefficient is not itself a direct causal effect. Compare models containing direct grandparental investments, transmitted unobserved traits, group persistence, and correlated measurement error that may yield similar associations. Seek additional data or exogenous variations that distinguish their quantitative contributions. Define any causal intervention separately, for example a feasible grandparental transfer holding specified parental resources fixed. Estimate decay of ancestry associations and transportability across family structures and historical periods. No first-order autoregression or universal persistence rate is imposed. Existing linked-data findings resolve some descriptive comparisons; the remaining research agenda concerns identifying underlying mechanisms under declared measurement and selection assumptions.

\textbf{Definitions and admissible scope.} A family line must specify whose ancestry is followed, how both parents and multiple grandparents are represented, and how migration, mortality, and unlinked records enter sampling. Let lifetime status \(Y_g\) have finite second moments and nonzero variance. The observational higher-ancestry test is conditional independence \(Y_g\mathrel{\perp}Y_{g-2}\mid(Y_{g-1},X)\), not a zero coefficient in an arbitrary linear regression. Measurement error is a joint distribution for repeated proxies conditional on true status; classical independence is an additional restriction, not a default. A direct grandparent effect compares a stated grandparent transfer or contact intervention at fixed parental inputs. Neither rejection of conditional independence nor slower-than-geometric correlations alone establishes that effect. Index observed generations by \(g=0,\ldots,G\), \(G\ge2\), with a fixed family-linking and ancestry-aggregation rule when a person has several ancestors. The first-order descriptive null is \(Y_g\mathrel{\perp}(Y_0,\ldots,Y_{g-2})\mid(Y_{g-1},X)\) on the declared common support. Identification of this latent-status null requires a measurement-error law, replicate indicators, validation data, or stated error bounds; the same test on proxies is a different object. Specify linkage-selection probabilities and cohort conditioning as inputs. A causal grandparent transfer contrast uses separately defined feasible transfer paths and their parental responses, with no causal interpretation assigned merely to rejection of the descriptive null.
\subsection*{Variants and assumption changes}
\begin{enumerate}
\item Measurement validation (source-motivated): add repeated independent status proxies or administrative lifetime records. Contrast latent first-order dependence with genuine higher-order dependence under an explicit error model.
\item Direct ancestral intervention (editorial): compare randomized grandparent transfers or exogenous contact opportunities under declared parental responses, measuring child outcomes. Test how this effect varies with co-residence and institutions instead of asserting a universal dynastic persistence rate.
\end{enumerate}
\subsection*{References and source audit}
\textbf{Support and access.} Full paper checked. The original p. 11 locator was inaccurate; mechanism discrimination is an explicit agenda, not a universal persistence law. \textbf{Locator:} IZA DP 10623, Section 2 final paragraph, printed p. 12; conclusion pp. 20--21. The 2017 working-paper date is separated from its later journal DOI.
\begin{enumerate}\raggedright
\item\hypertarget{definitions:source:30007646:1}{} Gary Solon. 2017. \emph{What Do We Know So Far about Multigenerational Mobility?}. IZA DP 10623, Section 2 final paragraph, printed p. 12; conclusion pp. 20--21.
\url{https://docs.iza.org/dp10623.pdf}.
\item\hypertarget{definitions:source:30007646:2}{} Gary Solon. 2018. \emph{What Do We Know So Far about Multigenerational Mobility?}. Publisher citation and abstract; The Economic Journal 128(612), F340--F352.
\url{https://onlinelibrary.wiley.com/doi/abs/10.1111/ecoj.12495}.
DOI: \url{https://doi.org/10.1111/ecoj.12495}.
\end{enumerate}

\subsection*{Common conventions: Source mathematical conventions}

These conventions apply to each entry unless explicitly replaced.

\textbf{Models and identification.} A primitive model \(m\) specifies all probability laws, feasible choices, information, equilibrium and measurement rules used by its target. The admissible class \(\mathcal M\) is fixed independently of outcomes. If \(P_O(m)\) is its observable probability law and \(\tau(m)\) the target, its identified set at observable law \(P\) is
\[
 \mathcal I_\tau(P)=\{\tau(m):m\in\mathcal M,\ P_O(m)=P\}.
\]
Point identification means that this set is a singleton. An empty set signals incompatibility of data and assumptions. Coordinate bounds need not describe the joint set. Bounds are sharp only if every retained target is attainable or arbitrarily approachable by compatible models. Finite bounds require explicit restrictions on unobserved quantities; unrestricted confounding, hidden wealth, extrapolation or arbitrary equilibrium selection can yield uninformative sets.

\textbf{Causal interventions.} A potential outcome \(Y_i(p)\) is the outcome under a complete policy regime \(p\), including timing, financing, information and market spillovers. The target population and units are fixed before treatment. With interference, use the complete assignment vector or a justified exposure mapping; individual no-interference notation alone cannot identify an economy-wide intervention. A difference of mean potential outcomes does not require the cross-world joint distribution. Conditioning on post-treatment migration, survival, enrollment or employment changes the estimand and needs separate justification.

\textbf{Statistical statements.} An estimator, confidence set, forecast or test specifies a sampling law, model class and loss/coverage criterion. Uniform \((1-\alpha)\)-coverage means \(\inf_{m\in\mathcal M}\Pr_m\{\tau(m)\in C_n\}\geq1-\alpha\), or an explicitly asymptotic version. The whole parameter space trivially has coverage, so informativeness requires a stated width, risk or power objective. A forecasting improvement is relative to a named benchmark, information set and evaluation population.

\textbf{Equilibrium and optimization.} An equilibrium comprises feasible plans, optimality, beliefs where needed, budget constraints and all market/resource clearing. The primitives or a stated benchmark define each correspondence \(\mathcal E(p;m)\). Nonemptiness is assumed or proved, never inferred just from compact policy sets. A selected-equilibrium target uses a declared selection rule; a robust target quantifies over all permitted equilibria. Use suprema or infima unless attainment is established. An \(\varepsilon\)-optimal policy is feasible and within \(\varepsilon\) of the finite optimum value. Report an empty feasible set separately.

\textbf{Welfare and accounting.} Normative utility, interpersonal normalization, social weights, numeraire, discounting and the use of public receipts are primitives. Behavioral utility need not coincide with the welfare criterion. Household welfare includes ownership income and rents once; adding profits again to compensating variation generally double-counts them. A monetary equivalent is defined by an explicit transfer and price convention, with monotonicity and range conditions. Average effects, mechanism contributions, transfers and welfare losses are different targets. Infinite sums require convergence and dynamic feasible plans require terminal or no-Ponzi conditions.

\textbf{Source versions.} A working-paper date, revision date and journal publication year are distinguished. A DOI for a journal version is not automatically the DOI of an earlier manuscript. Locators refer to the stated version and printed pages where specified. A metadata-only check does not verify a claim about a full-text conclusion. Every entry's source subsection narrows the provenance claim accordingly.
% END ECONOMICS STATEMENT
\end{document}
